% authority: science_draft
\documentclass[11pt]{article}

\usepackage[margin=1in]{geometry}
\usepackage{amsmath,amssymb,amsthm}
\usepackage{booktabs}
\usepackage{longtable}
\usepackage{enumitem}

\newcommand{\EqSrc}{\mathrm{EqSrc}}
\newcommand{\EqSrcT}{\mathrm{EqSrc}_{T}}
\newcommand{\Block}{\mathrm{Block}}
\newcommand{\NotEqSrc}{\mathrm{not\_EqSrc}}
\newcommand{\RetainH}{\mathrm{RetainH}}
\newcommand{\GenH}{\mathrm{GenH}}
\newcommand{\Fsrc}{\mathcal{F}_{\mathrm{src}}}
\newcommand{\Osrc}{\mathcal{O}_{\mathrm{src}}}
\newcommand{\Msrc}{\mathcal{M}_{\mathrm{src}}}
\newcommand{\Isrc}{\mathcal{I}_{\mathrm{src}}}
\newcommand{\Asrc}{\mathcal{A}_{\mathrm{src}}}
\newcommand{\Psrc}{\mathcal{P}_{\mathrm{src}}}
\newcommand{\Nsrc}{\mathcal{N}_{\mathrm{src}}}
\newcommand{\Csrc}{C_{\mathrm{src}}}
\newcommand{\NoTarget}{\mathsf{NoTarget}}

\newtheorem{countermodel}{Finite Countermodel}
\newtheorem{lemma}{Lemma}
\newtheorem{remark}{Boundary Remark}

\title{EqSrc Family-Closure Refuter Stress v1}
\author{Refuter Packet RT-20260707-023}
\date{2026-07-07}

\begin{document}
\maketitle

\section{Control Status}

This artifact implements v18 P3-T05 as a draft/control Refuter stress packet.
It stress-tests the P3-T02 conditional \(\EqSrcT\) family-closure theorem
candidate after the P3-T04 source-purity audit.

The stress result is:
\[
  \texttt{scoped\_obstruction}.
\]

\begin{verbatim}
stress_result: scoped_obstruction
refuter_classification: scoped_obstruction
loop_classification: scoped_obstruction
target_derivation_milestone: source_equivalence_eqsrc
milestone_burden: Stress the P3 family-closure result against closure removal,
  invariant weakening, primitive-dependence, and target-import attacks.
minimal_countermodel_available: true
next_route: P3-T06
general_EqSrc_adopted: false
RetainH_adopted: false
GenH_adopted: false
source_law_adopted: false
distance_to_gr_delta: no_distance_delta
\end{verbatim}

This result has no promotion authority. It does not discharge general
\(\EqSrc\), adopt \(\RetainH\), adopt \(\GenH\), adopt a source law, edit
canonical ontology, adopt \(\mathrm{MetricData}(E)\), expand
\(g_{\mathrm{eff}}\), derive matter coupling, derive Einstein equations,
promote a benchmark, issue a Gate Chair verdict, authorize external outreach,
or claim completed derivation.

\section{Inputs and Scope}

The object under stress is the P3-T02 conditional theorem candidate:
\begin{quote}
\texttt{research\_control/tasks/RT-20260707-020/artifacts/eqsrc\_family\_closure\_theorem\_or\_countermodel\_v1.tex}.
\end{quote}

The primitive-boundary input is:
\begin{quote}
\texttt{research\_control/tasks/RT-20260707-021/artifacts/retainh\_genh\_primitive\_boundary\_v1.tex}.
\end{quote}

The audit input is:
\begin{quote}
\texttt{research\_control/tasks/RT-20260707-022/artifacts/eqsrc\_family\_closure\_smuggling\_audit\_v1.tex}.
\end{quote}

The stress question is not whether the P3-T02 proof is source-pure as written.
P3-T04 already answered that in the affirmative. The stress question is
whether the conditional theorem candidate survives the ten P3-T05 attack modes
without becoming an adopted general \(\EqSrc\) theorem, a \(\RetainH\)
adoption, a \(\GenH\) adoption, or a target/process-authority overread.

\section{Stress Method}

For each stress mode, the Refuter removes or weakens one source-side
obligation and checks whether the positive family-closure conclusion still
follows. If the positive conclusion requires a missing closure witness, missing
ledger certificate, missing primitive, target premise, or process-authority
premise, the branch fails closed.

The P3-T02 theorem branch remains a conditional theorem candidate under
supplied H1--H7. The attacked overread is the stronger claim that the current
ontology has derived family closure generally or that the theorem candidate is
already adopted \(\EqSrc\).

\begin{verbatim}
stress_modes_checked:
  - remove_family_identity_closure
  - remove_inverse_closure
  - remove_composition_closure
  - weaken_invariant_ledger
  - expand_source_family_without_GenH
  - apply_H_retention_without_RetainH
  - replace_source_invariant_with_target_success
  - allow_missing_negative_controls
  - import_metric_or_detector_protocol
  - treat_theorem_candidate_as_adopted_EqSrc
\end{verbatim}

\section{Stress Mode Matrix}

\begin{longtable}{p{0.30\linewidth}p{0.20\linewidth}p{0.42\linewidth}}
\toprule
Mode & Result & Reason \\
\midrule
\texttt{remove\_family\_identity\_closure} & \texttt{obstruction} & Reflexivity is not derivable for an object whose source identity witness is absent from \(\Msrc\). \\
\texttt{remove\_inverse\_closure} & \texttt{finite\_countermodel} & The P3-T02 missing-inverse countermodel blocks symmetry with source data only. \\
\texttt{remove\_composition\_closure} & \texttt{finite\_countermodel} & The new three-object missing-composition countermodel below blocks transitivity. \\
\texttt{weaken\_invariant\_ledger} & \texttt{fails\_closed} & Without preservation of \(\Isrc,\Asrc,\Psrc\), \(\Csrc\) must return \(\Block\). \\
\texttt{expand\_source\_family\_without\_GenH} & \texttt{extension\_obstruction} & Generated objects cannot enter the theorem domain without a source-side \(\GenH\) rule. \\
\texttt{apply\_H\_retention\_without\_RetainH} & \texttt{extension\_obstruction} & H-retained records cannot be used as closure evidence without a source-side \(\RetainH\) rule. \\
\texttt{replace\_source\_invariant\_with\_target\_success} & \texttt{blocked} & Target success is excluded by \(\NoTarget\) and cannot replace \(\Isrc\) preservation. \\
\texttt{allow\_missing\_negative\_controls} & \texttt{fails\_closed} & Missing \(\Nsrc\) checks remove exclusion evidence and block positive comparison. \\
\texttt{import\_metric\_or\_detector\_protocol} & \texttt{blocked} & A metric or detector protocol would be a forbidden target/downstream premise. \\
\texttt{treat\_theorem\_candidate\_as\_adopted\_EqSrc} & \texttt{scoped\_obstruction} & Conditional theorem-candidate status does not equal adopted general \(\EqSrc\). \\
\bottomrule
\end{longtable}

\section{Missing-Inverse Countermodel Retained}

The P3-T02 missing-inverse countermodel remains valid stress evidence. It uses
\(\Osrc=\{A,B\}\), \(\Msrc=\{\mathrm{id}_{A},\mathrm{id}_{B},f:A\to B\}\),
and a source comparison rule with \(\Csrc(A,B)=\EqSrc\) but
\(\Csrc(B,A)=\Block\) because no source-side inverse \(B\to A\) is present.
Thus symmetry fails without inverse closure.

\section{Missing-Composition Countermodel}

\begin{countermodel}[Missing composition blocks transitivity]
Let
\[
  \Osrc=\{A,B,C\}
\]
and let
\[
  \Msrc=\{\mathrm{id}_{A},\mathrm{id}_{B},\mathrm{id}_{C},
  f:A\to B,\ g:B\to C\}.
\]
Assume every listed morphism is source-side, preserves the declared local
ledger, and satisfies \(\NoTarget\). Define
\[
  \Csrc(A,A)=\Csrc(B,B)=\Csrc(C,C)=\EqSrc,
  \quad \Csrc(A,B)=\EqSrc,\quad \Csrc(B,C)=\EqSrc,
\]
but
\[
  \Csrc(A,C)=\Block
\]
because no source-side composite \(h:A\to C\) is present in \(\Msrc\).
\end{countermodel}

Then \(A\,\EqSrcT\,B\) is witnessed by \(f\), and \(B\,\EqSrcT\,C\) is
witnessed by \(g\), but \(A\,\EqSrcT\,C\) is blocked. The construction uses
only source objects, source morphisms, a source comparison rule, ledger
preservation, and the no-target guard. Therefore transitivity is not derivable
without the composition-closure hypothesis.

\section{Ledger-Weakening Witness}

\begin{lemma}[Ledger weakening is fail-closed]
If a witness \(w:X\to Y\) no longer preserves the required entries of
\(\Isrc,\Asrc,\Psrc\), and no replacement source certificate is supplied, then
the source-only comparison rule cannot return \(\EqSrc\) for \(X,Y\) without
overreading the ledger.
\end{lemma}

\textbf{Reasoning.} The P3-T02 relation requires explicit preservation of the
invariant ledger, annotations, and forbidden-channel/proxy-edge incidence.
Weakening that ledger removes the certificate that the comparison is between
the same source-side content. The fail-closed comparison is
\(\Csrc(X,Y)=\Block\), not target repair and not process-authority repair.

\section{RetainH and GenH Extension Pressure}

P3-T03 classified \(\RetainH\) and \(\GenH\) as
\texttt{not\_required\_here} for the already declared closed-family theorem.
P3-T05 preserves that result.

The extension attacks have a narrower implication:
\begin{itemize}[leftmargin=*]
  \item \texttt{apply\_H\_retention\_without\_RetainH} blocks using
  H-retained records as theorem evidence until a source-side retention rule is
  supplied.
  \item \texttt{expand\_source\_family\_without\_GenH} blocks using generated
  objects or morphisms as theorem-domain members until a source-side generation
  rule is supplied.
\end{itemize}

These are scoped extension obstructions. They do not refute the closed-family
conditional theorem under H1--H7, and they do not adopt \(\RetainH\) or
\(\GenH\).

\section{Target and Process Authority Attacks}

The target-import attacks fail closed. Replacing source invariants with target
success, importing a target metric, importing a detector protocol, or using
proper-time calibration would violate \(\NoTarget\). Such premises cannot
repair a missing source inverse, missing composition, weak ledger, missing
\(\RetainH\), missing \(\GenH\), or missing negative-control evidence.

The process-authority attacks also fail closed. Registry rows, validator
passes, role identity, handoff status, commit state, and generated derivative
status are provenance and operational receipts only. They do not make source
objects, morphisms, composites, ledger certificates, \(\RetainH\), or
\(\GenH\) exist.

\section{Refuter Obstruction Record}

\begin{verbatim}
refuter_obstruction_record:
  obstruction_id: "OB-V18-P3T05-EQSRC-FAMILY-CLOSURE-001"
  target_claim: "P3-T02 conditional EqSrc_T family-closure theorem candidate
    read as adopted general EqSrc discharge"
  target_milestone: "source_equivalence_eqsrc"
  failed_premise: "H1-H7 closure and ledger hypotheses are supplied
    conditionally and are not derived by current ontology for closure-removed
    or extension-expanded cases"
  minimal_countermodel_available: true
  countermodel_path: "research_control/tasks/RT-20260707-023/artifacts/eqsrc_family_closure_refuter_stress_v1.tex"
  countermodel_scope: "finite_local"
  certificate_gap: "missing source-side inverse composition or
    ledger-preservation certificate under the stressed modes"
  source_extension_repair_possible: true
  global_no_go_claim_authorized: false
  future_source_extension_impossibility_authorized: false
  freeze_criteria_status:
    freeze_decision: "not_frozen"
    decision_reason: "This is the first P3-T05 stress of this exact
      family-closure burden and P3-T06 selector remains mandated."
    next_allowed_route: "P3-T06"
  route_cycle_control:
    cycle_family: "v18_p3_eqsrc_family_closure"
    current_cycle_step: "refuter_stress"
    prior_related_tasks:
      - "RT-20260707-020"
      - "RT-20260707-021"
      - "RT-20260707-022"
    cycle_risk: "medium"
    orbit_avoidance_reason: "The packet adds a concrete
      missing-composition finite countermodel and routes to selector rather
      than repeating theorem construction."
    next_role_consequence: "theoretical-continuation-selector@0.1.0"
  forbidden_conclusions:
    - "No general EqSrc discharge follows."
    - "No RetainH adoption follows."
    - "No GenH adoption follows."
    - "No source-law adoption follows."
    - "No MetricData(E) adoption follows."
    - "No g_eff scope expansion follows."
    - "No matter-coupling derivation follows."
    - "No Einstein-equation derivation follows."
    - "No benchmark promotion follows."
    - "No Gate Chair verdict follows."
    - "No completed derivation follows."
    - "No program-wide no-go conclusion follows."
    - "No future source-extension impossibility follows."
\end{verbatim}

\section{Stress Result}

The stress result is \(\texttt{scoped\_obstruction}\).

The conditional theorem candidate survives exactly as a theorem candidate under
supplied H1--H7. The overread fails: current ontology has not derived those
hypotheses generally, and closure-removed or extension-expanded branches
produce finite/local source-side obstructions. The minimal countermodel status
therefore survives, but only inside the named finite closure-removal scope.

\begin{verbatim}
stress_result_one_of:
  - survives_as_theorem_candidate
  - minimal_countermodel_survives
  - repair_required
  - scoped_obstruction
  - freeze_route
selected_stress_result: scoped_obstruction
minimal_countermodel_survives: true
repair_required: false
freeze_route: false
same_milestone_continuation_status: blocked_adoption_open_continuation
\end{verbatim}

\section{Distance-to-GR Effect}

\begin{longtable}{p{0.30\linewidth}p{0.56\linewidth}}
\toprule
Burden & Status after this packet \\
\midrule
Source ontology primitives & Unchanged; no canonical ontology edit. \\
Source equivalence EqSrc & Family-closure overread stress returns scoped obstruction; no general EqSrc discharge. \\
RetainH & Not adopted; extension pressure remains candidate-definition-needed. \\
GenH & Not adopted; generated-family extension pressure remains candidate-definition-needed. \\
ObsLoc\_lc & Unchanged. \\
Resp\_lc & Unchanged. \\
M\_src & Unchanged. \\
g\_eff & Unchanged; no MetricData(E) or metric-scope expansion. \\
matter coupling & Blocked and unchanged; no coupling derivation or adoption. \\
Einstein equations & Blocked and unchanged; no field-equation derivation. \\
finite-variation robustness & Unchanged; no global robustness theorem. \\
benchmark promotion & Blocked; no benchmark evidence or promotion. \\
Gate Chair status & Unchanged; no Gate Chair verdict. \\
current route freeze or hard-fail status & Not frozen; route continues to P3-T06 selector. \\
\bottomrule
\end{longtable}

\begin{verbatim}
distance_to_gr_delta:
  changed: false
  effect: no_distance_delta
  ledger_row_updated: false
  reason: P3-T05 records a scoped Refuter obstruction and finite local
    closure-removal countermodels only. It does not adopt EqSrc RetainH GenH
    or any downstream physics claim.
\end{verbatim}

\section{Forbidden Conclusions}

This Refuter stress result must not be read as any of the following:
\begin{itemize}[leftmargin=*]
  \item general \(\EqSrc\) discharge;
  \item \(\RetainH\) adoption;
  \item \(\GenH\) adoption;
  \item source-law adoption;
  \item canonical ontology edit;
  \item target metric or detector-protocol import;
  \item \(\mathrm{MetricData}(E)\) adoption;
  \item \(g_{\mathrm{eff}}\) scope expansion;
  \item source detector/readout semantics adoption;
  \item matter-semantics adoption;
  \item coupling-law adoption;
  \item matter-coupling derivation or adoption;
  \item stress-energy semantics, a stress-energy tensor, or a matter action;
  \item Einstein-equation derivation;
  \item benchmark promotion;
  \item Gate Chair verdict;
  \item completed derivation;
  \item program-wide rejection claim; or
  \item future source-extension impossibility.
\end{itemize}

\section{Next Route}

The next route is:
\[
  \texttt{P3-T06}.
\]

P3-T06 must select the post-family-closure continuation route. The default
route remains P4-T01 unless the selector determines that repair, freeze review,
\(\RetainH\), \(\GenH\), source detector/readout semantics, or external review
preparation is more appropriate under the v18 plan.

\section{Source Materials}

The AEther-Flow Research Project. (2026, July 7). \textit{EqSrc
family-closure theorem or countermodel v1} [Internal science draft].
\texttt{research\_control/tasks/RT-20260707-020/artifacts/eqsrc\_family\_closure\_theorem\_or\_countermodel\_v1.tex}

The AEther-Flow Research Project. (2026, July 7). \textit{RetainH and GenH
primitive-boundary extraction v1} [Internal science draft].
\texttt{research\_control/tasks/RT-20260707-021/artifacts/retainh\_genh\_primitive\_boundary\_v1.tex}

The AEther-Flow Research Project. (2026, July 7). \textit{EqSrc
family-closure smuggling audit v1} [Internal science draft].
\texttt{research\_control/tasks/RT-20260707-022/artifacts/eqsrc\_family\_closure\_smuggling\_audit\_v1.tex}

The AEther-Flow Research Project. (2026, July 7). \textit{Recommendations
implementation plan continue task v18} [Internal implementation plan].
\texttt{implementations\_plans/recommendations\_implementation\_plan\_continue\_task-v18.md}

The AEther-Flow Research Project. (2026, July 2). \textit{Refuter obstruction
schema v1} [Internal project-control note].
\texttt{research\_control/design/refuter\_obstruction\_schema\_v1.md}

\end{document}
